%!TEX root = ../main.tex

% ============================================================
%  CAPÍTULO 7: IMPACTO ESPERADO
% ============================================================

\chapter{Impacto esperado}\label{ch:impacto}

% ──────────────────────────────────────────────────────────────
\section{Impacto académico}
% ──────────────────────────────────────────────────────────────

Esta investigación aspira a realizar una aportación fundacional a la literatura de ciencia de datos políticos y estudios electorales cuantitativos en América Latina. A nivel epistemológico, este trabajo documenta la primera metodología estructurada de \textit{nowcasting} electoral elección-agnóstica e independiente de encuestas conceptualizada e implementada para el entorno político altamente fragmentado de Colombia. 

Metodológicamente, la formalización matemática del \textbf{modelo de crecimiento Weibull} constituye un artefacto científico de exportación. Al proporcionar un marco matemático capaz de reconstruir la evolución temporal de la tracción digital sin depender de registros en tiempo real, el estudio libera a futuras investigaciones académicas de la obligación logística de desplegar sensores de recopilación con meses de antelación al evento electoral de interés. 

Adicionalmente, se efectúa una contribución transversal al debate global sobre la viabilidad de la predicción electoral mediante redes sociales, delineando claramente las fronteras de escala empírica: delimitando la inviabilidad estructural de las arquitecturas de regresión volumétrica en escenarios de escasez dimensional (muestras menores a 200 casos), al tiempo que valida matemáticamente la \textit{dominancia relativa} como el estimador heurístico más robusto en elecciones subnacionales.

% ──────────────────────────────────────────────────────────────
\section{Impacto práctico}
% ──────────────────────────────────────────────────────────────

En una dimensión operativa, el principal aporte aplicativo es la materialización de un sistema orquestado de pronóstico y monitoreo de extremado bajo costo. Al demostrarse la reproducibilidad absoluta del ecosistema metodológico empleando herramientas comerciales de extracción \textit{off-the-shelf} (como \textit{Apify} o integradores \textit{no-code} como \textit{Make.com}) y prescindiendo por completo de servidores y arquitecturas universitarias dedicadas, la investigación democratiza el análisis de datos políticos complejos.

Este avance abre la puerta para que medios de comunicación regionales, laboratorios de veeduría ciudadana y pequeñas firmas de consultoría estratégica desplieguen capacidades de inteligencia electoral en tiempo real. En el mediano plazo, estas técnicas se proyectan como la alternativa natural y costo-eficiente frente a los oligopolios demoscópicos, especialmente en otras naciones de Latinoamérica donde la carencia de encuestas sistemáticas o la desconfianza pública en sus resultados exigen métodos de estimación paralelos transparentes y auditables.

% ──────────────────────────────────────────────────────────────
\section{Limitaciones y trabajo futuro}
% ──────────────────────────────────────────────────────────────

La honestidad intelectual obliga a delimitar las fallas orgánicas y los confines de aplicabilidad del estudio. En primer lugar, la calibración de la matriz temporal Weibull, si bien fue exhaustiva a nivel micro de publicaciones, descansó sobre el comportamiento macro de una única elección (Costa Rica 2026). Asumir universalidad conductual exige una validación prospectiva replicando la medición longitudinal sobre certámenes demográficamente disímiles para confirmar la estabilidad inter-elecciones de los parámetros de forma y escala. 

En segundo lugar, por determinaciones corporativas restrictivas de la compañía Meta, Instagram quedó excluido del \textit{dataset} colombiano; esta omisión representa una pérdida incalculable de la fotografía generacional y el consumo audiovisual de una gran franja del electorado joven. Asimismo, tanto el modelo de dominancia heurística como los ensambles algorítmicos probaron ser estructuralmente vulnerables a perfiles políticos con audiencias atípicas (como lo demostró la exclusión forzosa del candidato \textit{influencer} Gustavo Bolívar), revelando que las masas críticas interdepartamentales contaminan la medición estrictamente local. Todo esto sin ignorar que el conteo volumétrico peca de ceguera ante la calidad y veracidad de la interacción: las arquitecturas modeladas no están equipadas actualmente para ponderar el efecto adverso o positivo de interacciones manufacturadas artificialmente (granjas de \textit{bots}, perfiles inorgánicos o astroturfing).

Por lo tanto, la agenda de **trabajo futuro** a desprender de esta investigación es rica y expansiva. Se priorizará escalar el conjunto de datos integrando eventos asimétricos (elecciones legislativas departamentales) y otras geografías de la región andina; investigar el desempeño matricial de agregar ventanas temporales de memoria (\textit{lagged variables}); desarrollar clasificadores semánticos que depuren la señal descartando la actividad inorgánica artificial; e incorporar en el plano explicativo covariables estructuradas \textit{offline} (indicadores económicos territoriales y clima de opinión en la prensa tradicional). Inmediatamente, la prioridad empírica definitiva será la confrontación longitudinal ciega del clasificador binario entrenado contra los resultados de las próximas elecciones territoriales en Colombia del año 2027.
